\chapter{Konstruktionen und Ununterscheidbarkeit}
\label{b11guide:ch:constructions}

\section{Einschränken und Zusammenfügen}

\begin{definition}[Eingeschränkte Äquivalenzrelation]
Für $E\in\EqSet(A)$ und $B\subseteq A$ sei
\[
 E|_B=E\cap(B\times B).
\]
\end{definition}

\begin{theorem}[Einschränkung und Klassen]
\label{b11guide:thm:restriction}
Für $E\in\EqSet(A)$ und $B\subseteq A$ gilt
\[
\begin{aligned}
 E|_B&\in\EqSet(B),\\
 [x]_{E|_B}&=[x]_E\cap B\qquad(x\in B),\\
 B/(E|_B)&=\{C\cap B\mid C\in A/E,\ C\cap B\ne\varnothing\}.
\end{aligned}
\]
Ist $B$ zusätzlich $E$-gesättigt, so gilt
\[
 B/(E|_B)=\{C\in A/E\mid C\subseteq B\}.
\]
\end{theorem}

\begin{theorem}[Erweiterung durch Einzelklassen]
Seien $B\subseteq A$ und $E\in\EqSet(B)$. Dann ist
\[
 E^+=E\cup\Delta_{A\setminus B}
\]
eine Äquivalenzrelation auf $A$, und
\[
 E^+|_B=E,\qquad
 A/E^+=(B/E)\cup\{\{x\}\mid x\in A\setminus B\}.
\]
Für jede Äquivalenzrelation $F$ auf $A$ mit $E\subseteq F$
gilt $E^+\subseteq F$.
\end{theorem}

\begin{theorem}[Disjunkte Vereinigung]
Sei $(A_i)_{i\in I}$ eine Familie paarweise disjunkter Mengen,
und für jedes $i\in I$ sei eine Relation $E_i\in\EqSet(A_i)$
gegeben. Für
\[
 A=\bigcup_{i\in I}A_i,\qquad E=\bigcup_{i\in I}E_i
\]
gilt
\[
 E\in\EqSet(A),\qquad
 [x]_E=[x]_{E_i}\ (x\in A_i),\qquad
 A/E=\bigcup_{i\in I}(A_i/E_i).
\]
Die Aussage umfaßt $I=\varnothing$.
\end{theorem}

\section{Produkte und Durchschnitte}

\begin{definition}[Produkt zweier Äquivalenzrelationen]
Für $E\in\EqSet(A)$ und $F\in\EqSet(B)$ sei die Relation
$E\boxtimes F$ auf $A\times B$ durch
\[
 (a,b)\sim_{E\boxtimes F}(a',b')
 \quad\Longleftrightarrow\quad
 a\sim_E a'\ \land\ b\sim_F b'
\]
definiert.
\end{definition}

\begin{theorem}[Produktklassen]
\label{b11guide:thm:product}
Für $E\in\EqSet(A)$ und $F\in\EqSet(B)$ gilt
\[
\begin{aligned}
 E\boxtimes F&\in\EqSet(A\times B),\\
 [(a,b)]_{E\boxtimes F}&=[a]_E\times[b]_F
                       \qquad((a,b)\in A\times B),\\
 (A\times B)/(E\boxtimes F)
     &=\{C\times D\mid C\in A/E,\ D\in B/F\}.
\end{aligned}
\]
Ist einer der Träger leer, sind beide Seiten der letzten
Gleichheit leer.
\end{theorem}

\begin{theorem}[Beliebiger Durchschnitt]
\label{b11guide:thm:intersection}
Sei $\mathcal E\subseteq\EqSet(A)$. Die Relation
\[
 H=\{(x,y)\in A\times A\mid
                   \forall E\in\mathcal E\;(x\sim_E y)\}
\]
ist eine Äquivalenzrelation auf $A$. Für $x\in A$ gilt
\[
 [x]_H=\{y\in A\mid\forall E\in\mathcal E\;(y\in[x]_E)\}.
\]
Bei $\mathcal E\ne\varnothing$ sind dies
$H=\bigcap\mathcal E$ und $[x]_H=\bigcap_{E\in\mathcal E}[x]_E$.
Für $\mathcal E=\varnothing$ gilt $H=\nabla_A$ und $[x]_H=A$.
\end{theorem}

\begin{theorem}[Gemeinsame Klassenverfeinerung zweier Relationen]
Für $E,F\in\EqSet(A)$ gilt
\[
\begin{aligned}
 [x]_{E\cap F}&=[x]_E\cap[x]_F\qquad(x\in A),\\
 A/(E\cap F)&=\{C\cap D\mid C\in A/E,\ D\in A/F,\
                                      C\cap D\ne\varnothing\}.
\end{aligned}
\]
\end{theorem}

\section{Ununterscheidbarkeit durch Mengen}

\begin{definition}[Ununterscheidbarkeitsrelation]
\label{b11guide:def:indiscernibility}
Für $\mathcal M\subseteq\mathcal P(A)$ sei $E_{\mathcal M}^A$
die Relation auf $A$ mit
\[
 x\sim_{E_{\mathcal M}^A}y
 \quad\Longleftrightarrow\quad
 \forall M\in\mathcal M\;(x\in M\Longleftrightarrow y\in M).
\]
Für $A=\bigcup\mathcal M$ ist dies die
Ununterscheidbarkeitsrelation aus Band~11.
\end{definition}

\begin{theorem}[Ununterscheidbarkeit und ihre Klassen]
Für $\mathcal M\subseteq\mathcal P(A)$ gilt
$E_{\mathcal M}^A\in\EqSet(A)$. Alle Mengen $M\in\mathcal M$
sind bezüglich dieser Relation gesättigt.
Für $x\in A$ gilt
\[
 [x]_{E_{\mathcal M}^A}
 =
 \left(\bigcap_{\substack{M\in\mathcal M\\x\in M}}M\right)
 \cap
 \left(\bigcap_{\substack{M\in\mathcal M\\x\notin M}}(A\setminus M)\right).
\]
Leere Durchschnitte werden relativ zu $A$ als $A$ verstanden.
Insbesondere gilt $E_{\varnothing}^A=\nabla_A$.
\end{theorem}

\begin{theorem}[Darstellung jeder Äquivalenz durch Mengenzugehörigkeit]
\label{b11guide:thm:recover}
Für $E\in\EqSet(A)$ gilt
\[
 E_{A/E}^A=E,\qquad E_{\SatSets_E(A)}^A=E.
\]
Zwei Elemente sind also genau dann $E$-äquivalent, wenn sie
in sämtlichen $E$-Klassen dieselbe Mitgliedschaft haben;
gleichwertig kann man alle gesättigten Mengen verwenden.
Für $\mathcal M\subseteq\mathcal N\subseteq\mathcal P(A)$ gilt
\[
 E_{\mathcal N}^A\subseteq E_{\mathcal M}^A.
\]
Zusätzliche Mitgliedschaftsbedingungen können Klassen teilen,
aber keine vorher getrennten Klassen vereinigen.
\end{theorem}

